\begin{section}{experience}

  \begin{work}
    {Graduate Assistant}
    {Institute for Systems Research}
    {University of Maryland}
    {Aug 2023 - Present}

    % \item Served as a Teaching Assistant for a Robotics Laboratory Course:
    %       ENEE467 in Fall '23 semester, worked with students from
    %       diverse engineering backgrounds. Currently improving course material for the upcoming
    %       semester.

    % \item Assisted students with setting up GPIOs, SysTick timer, General
    %       purpose timer and PWM signal generation on TI-MSP432 Launchpad board.

    % \item Guided students in interdisciplinary lab exercises which included
    %       writing a roscpp node to draw shapes with a real UR3e robotic arm and
    %       developing a line follower program for the TI Robotic Systems Learning
    %       Kit.

    \item Developing an application to track and evaluate line following robots visually using fiducial
          markers. Utilized modern C++ (C++11) and OOP principles, OpenCV and Matplot++ for
          implementation. Results serve as a rubric for grading lab exercises.

    % \item Developing a program to track and evaluate line following robots visually using fiducial
    %       markers. Utilized modern C++ (C++11) and OOP principles, documented the software using
    %       Sphinx. Results serve as a rubric for grading lab exercises.

    \item Ported lab exercises from ROS1 to ROS2 by creating a custom Docker image and
          integrating it with VSCode to improve the development environment for students.

    \item Delivered a hotfix for a C++ wrapper library based on MoveIt!
          API used during the course. This significantly improved execution
          success rate of student-written code when using the Cartesian
          motion planner to move the robotic arm.

    \item Migrated the documentation for lab exercises from
          GitHub Wiki to Sphinx and hosted it on Read The Docs, resulting in
          centralized documentation for improved organization and maintenance.

    % \item Assisted students with debugging C++ and Embedded C code in lab
		% 			exercises.

    % \item Supported students in debugging Embedded C code by identifying
    %       incorrect register values and program flow errors.

  \end{work}

  \begin{work}
    {Robotics Engineer Intern}
    {iTrontik Smart Systems}
    {Pune}
    {May 2021 - Apr 2022}

    \item Developed hardware for a nano class (less than 250 grams) quadcopter drone for aerial
          photography applications.

    \item Sized the brushless motors and propellers for the drone, selected suitable off-the-shelf
          parts that achieved 90 percent of the target hover flight time on average in outdoor testing,
          with a 12 percent excess in total weight.

    % \item Designed a sturdy 3D-printed quadcopter frame for prototyping using
		% 			Fusion 360 in three iterations. The frame was able to sustain multiple
		% 			hard and crash landings during testing.

    \item Analyzed vibrations in the frame design using flight log data and modal analysis in Fusion
          360, and made improvements to the design which resulted in a 66 percent reduction of
          vibrations, making the drone flyable.

%     \item Created circuit schematic and 2-layer PCB design for a USB-C charger
%           board that integrates into the drone battery pack. This reduced the
%           need for an external charger.
		% \item Ideated and developed a USB-C charger board which integrates into the
		% 			battery pack to reduce the need to ship an external charger with the
		% 			drone. Designed a 2-layer PCB in KiCAD and worked with a PCB
		% 			fab house to produce the first prototype.

    \item Ideated and prototyped a USB-C charging solution for the drone battery pack, enabling
          the use of existing smartphone chargers instead of a dedicated external charger.

  \end{work}

  % \begin{work}
  %   {Software Team Member}
  %   {Terraformers URC}
  %   {}
  %   {Jan 2024 - May 2024}

  %   \item Written a program in C++ to evaluate the workspace of the arm using Monte Carlo simulation.
  %         The program is currently being used to optimize the DH parameters of the robotic arm for
  %         maximum reachability in the regions of interest.

  %   \item Built a URDF description package of the rover for simulating it using the ROS ecosystem.

  %   \item Maintained the projects and their documentation on GitHub repositories.

  % \end{work}
\end{section}
